\chapter{Селектор коэффициента центральной щели синглета и триплета}
\label{ch:v8-baryon-c0-singlet-triplet-central-gap-coefficient-selector-gate}

Направление $Q=P_{\mathbf3}-3I_4/4$ уже условно единственно. Проверим,
может ли естественная нормировка этого оператора выбрать физический
коэффициент $\lambda$ в $h=\varepsilon I_4+\lambda Q$.

\section{Шесть неэквивалентных нормировок}

Точные инварианты выбранного следово-центрированного представителя равны
\begin{equation}
 \Tr Q^2=\frac34,
 \qquad \lVert Q\rVert_{\rm op}=\frac34,
 \qquad \lVert Q\rVert_1=\frac32,
 \qquad \frac14\Tr Q^2=\frac3{16}.
 \label{eq:v8-baryon-c0-central-gap-coefficient-q-invariants}
\end{equation}
Проекторная и grading/Casimir-нормировки дают соответственно
\begin{equation}
 |\lambda|_{P}=1,
 \qquad
 |\lambda|_{\Gamma,\mathcal C_2}=2.
 \label{eq:v8-baryon-c0-central-gap-projector-grading-normalizations}
\end{equation}
Единичные гильбертово-шмидтова и операторная нормы дают
\begin{equation}
 |\lambda|_{\rm HS}=\frac2{\sqrt3},
 \qquad
 |\lambda|_{\rm op}=\frac43.
 \label{eq:v8-baryon-c0-central-gap-hs-operator-normalizations}
\end{equation}
Единичные следовая норма и дисперсия относительно $I_4/4$ дают ещё два
значения:
\begin{equation}
 |\lambda|_1=\frac23,
 \qquad
 |\lambda|_{\rm var}=\frac4{\sqrt3}.
 \label{eq:v8-baryon-c0-central-gap-trace-variance-normalizations}
\end{equation}
Таким образом,
\begin{equation}
 \left\{1,2,\frac2{\sqrt3},\frac43,\frac23,\frac4{\sqrt3}\right\}
 \quad\text{содержит шесть различных положительных значений}.
 \label{eq:v8-baryon-c0-central-gap-six-distinct-normalizations}
\end{equation}
Ни одно из них не имеет привилегии без указания, какая именно норма входит
в родительское действие.

Более того, норма нецентрированного гамильтониана зависит от произвольного
нуля энергии:
\begin{equation}
 \Tr(Q+bI_4)^2=4b^2+\frac34.
 \label{eq:v8-baryon-c0-central-gap-uncentered-norm-shift-dependence}
\end{equation}
Поэтому нормировочный контракт обязан сначала выбрать центрирование; здесь
оно уже использует обычный счётный след.

\section{Тепловая масштабная орбита}

Равновесный вес видит только
\begin{equation}
 \theta=\beta\lambda,
 \qquad
 (\beta,\lambda)\longmapsto
 \left(\frac\beta a,a\lambda\right),
 \qquad a>0.
 \label{eq:v8-baryon-c0-central-gap-coefficient-thermal-scale-orbit}
\end{equation}
Следовательно, даже выбранная безразмерная нормировка $Q$ не превращается в
физическую энергию без температурного или временного якоря.

\section{Чётный потенциал выбирает коэффициенты, а не число}

Минимальный устойчивый чётный потенциал имеет вид
\begin{equation}
 V_{a,b}(\lambda)=\frac a2\lambda^2+\frac b4\lambda^4,
 \qquad b>0.
 \label{eq:v8-baryon-c0-central-gap-even-quartic-potential}
\end{equation}
Его уравнение стационарности равно
\begin{equation}
 V'_{a,b}(\lambda)=\lambda(a+b\lambda^2)=0.
 \label{eq:v8-baryon-c0-central-gap-even-potential-stationarity}
\end{equation}
Два точных устойчивых выбора коэффициентов дают разные модули:
\begin{equation}
 (a,b)=(-1,1)\Longrightarrow |\lambda_*|=1,
 \qquad
 (a,b)=(-4,1)\Longrightarrow |\lambda_*|=2.
 \label{eq:v8-baryon-c0-central-gap-quartic-magnitude-witnesses}
\end{equation}
В обоих случаях минимумы попарно вырождены,
\begin{equation}
 V_{a,b}(\lambda)=V_{a,b}(-\lambda),
 \label{eq:v8-baryon-c0-central-gap-even-potential-sign-degeneracy}
\end{equation}
так что чётный parent не выбирает знак, а модуль определяется свободным
отношением $-a/b$.

\section{Энтропия и KMS не дают коэффициент}

Максимум центральной энтропии соответствует
\begin{equation}
 p=\frac14
 \quad\Longleftrightarrow\quad
 \theta=0
 \quad\Longleftrightarrow\quad
 \lambda=0\ \text{при фиксированном }\beta.
 \label{eq:v8-baryon-c0-central-gap-maximum-entropy-zero-gap}
\end{equation}
Это выбирает отсутствие щели, а не её ненулевой коэффициент. Равный полный
вес секторов и KMS-отношение дают лишь целевые уравнения
\begin{equation}
 p=\frac12\Longrightarrow\beta\lambda=\log3,
 \qquad
 r_{\rm obs}\Longrightarrow\beta\lambda=-\log r_{\rm obs}.
 \label{eq:v8-baryon-c0-central-gap-target-dependent-equations}
\end{equation}
Первое использует выбранную секторную нормировку, второе --- наблюдаемое
отношение, поэтому оба являются калибровками, а не слепыми селекторами.

Итоговый реестр равен
\begin{equation}
 N_{\rm coefficient\ selector}=0/8:
 \quad\{P,\Gamma,\mathcal C_2,\lVert\cdot\rVert_{\rm HS},
 \lVert\cdot\rVert_{\rm op},\operatorname{Var},V_{a,b},S/KMS\}.
 \label{eq:v8-baryon-c0-central-gap-coefficient-selector-ledger}
\end{equation}
Три проектные записи формы совпадают, но их численные нормировки и пять
дополнительных критериев не дают одного физического $\lambda$.

Минимальная новая родительская структура должна содержать одновременно
\begin{equation}
 V_{\rm src}(\lambda)=\frac{m^2}{2}\lambda^2-j\lambda,
 \qquad m^2>0,
 \qquad \lambda_*=\frac{j}{m^2},
 \label{eq:v8-baryon-c0-central-gap-minimal-source-preview}
\end{equation}
то есть чётную жёсткость и нечётный центральный источник. Следующий гейт
проверит минимальность и допустимость этой архитектуры, не приписывая ей
заранее происхождение.

\begin{theorem}[Нормировочная неединственность коэффициента щели]
Один и тот же центральный оператор $Q$ допускает как минимум шесть точных
естественных нормировок с различными значениями $|\lambda|$. Чётный
устойчивый потенциал выбирает модуль только через свободные коэффициенты и
не снимает знак. Поэтому физический коэффициент щели не следует из выбора
нормы представителя.
\end{theorem}

\begin{nogocriterion}[Запрет превращения единичной нормы в энергию]
Условие $\lVert\lambda Q\rVert=1$ является координатной нормировкой до тех
пор, пока та же норма с фиксированным коэффициентом не выведена из общего
действия. Нельзя отождествлять полученное безразмерное число с энергией,
температурой или скоростью релаксации.
\end{nogocriterion}

\begin{protocol}\begin{enumerate}
\item проверенных естественных нормировок: \textbf{$6$};
\item различных положительных модулей: \textbf{$6$};
\item норма без центрирования инвариантна к сдвигу: \textbf{нет};
\item равновесие видит отдельно $\beta$ и $\lambda$: \textbf{нет};
\item чётный потенциал выбирает модуль без коэффициентов: \textbf{нет};
\item чётный потенциал выбирает знак: \textbf{нет};
\item максимум энтропии: \textbf{$\lambda=0$};
\item KMS без target выбирает $\lambda$: \textbf{нет};
\item селекторов коэффициента: \textbf{$0/8$};
\item минимальная недостающая структура: \textbf{нечётный источник плюс жёсткость};
\item следующий гейт: \textbf{минимальная source-parent архитектура}.
\end{enumerate}\end{protocol}

Машинный сертификат сохранён в
\path{s2t_v8_baryon_c0_singlet_triplet_central_gap_coefficient_selector_gate_results.json}.